\scanpage{372}
\mspage{96}
plonger canoniquement dans un $\mathrm{Hot}_W$ convenable, où $W$ \unsure{serait} un localiseur fondamental strictement plus fin que $\mathbb W_\infty$, p. ex. (qui sait ?) $\widetilde{\mathbb W}_\infty$. Le premier travail qu'on aurait envie de faire, c'est de construire un \unsure{tel|telle} $W$ ad-hoc, pour que l'invariant de Whitehead puisse s'écrire comme un foncteur sur $\mathrm{Hot}_{\bullet W}$. Pour commencer, il faudrait déjà s'assurer que cet invariant \deleted{peut être} \unsure{en|est} \unsure{nomme|somme|vraiment} un invariant des types d'homotopie, i.e qu'il s'interprète comme un foncteur sur $\unsure{\mathrm{Hot}_h|\mathrm{Hot}_2|\mathrm{Hot}_\bullet|\mathrm{Hot}}\mathrm{Cat}\simeq\mathrm{Cat}(\underbrace{\mathrm{homotop.}}_{\substack{\text{ens. des}\\\text{homotopismes}}})^{-1}$.

Ce travail \unsure{suggèrera|suggérera|suggerera} peut-être alors la définition des ``invariants de Whitehead supérieurs'', (impliquant les $\pi_i$, non seulement \unsure{les|le} $\pi_1$), lesquels pourraient peut-être servir (avec celui de dim $1$) à décrire le localiseur hypothétique $W$ (``de Whitehead'').
% archive sheet number: 371
